\chapter{Homogenization and Interpolation}\label{chap:homogenisation-and-interpolation}

This chapter shall be devoted to looking at some structural results on arithmetic circuits. 
This would help us understand the relevance of shallow circuits in the context of proving lower bounds for arithmetic circuits of arbitrary depth. 

\section{Homogenization}\label{sec:homogenization}

Suppose we have an $n$-variate degree $d$ polynomial computed by an arithmetic circuit $C$. 
How large can the degree of intermediate computations be? 
Potentially, intermediate computations can involve very high degree terms which somehow cancel each other at the root. 
However, the following lemma of Strassen shows that we may assume without much loss of generality that arithmetic circuits never compute polynomials of degree more than the output. 

\begin{definition}[Homogeneous circuits]
A circuit $C$ is said to be \emph{homogeneous} if every gate in the circuit computes a homogeneous polynomial. 
\end{definition}

\begin{lemma}[Homogenization]\label{lem:homogenization}
Let $f$ be an $n$-variate degree $d$ polynomial computed by a circuit $C$ of size $s$. 
Then, for every $0\leq i \leq d$, there is a \emph{homogeneous arithmetic circuit} $C_i'$, of size at most $O(sd^2)$, that computes the degree $i$ homogeneous polynomial in $f$ (the sum of all degree $i$ monomials in $f$).  
\end{lemma}
\begin{proof}
Assume without loss of generality that the circuit $C$ has all gates with fan-in at most $2$. 
For every gate $g\in C$, define $(d+1)$ gates $g^{(0)},\dots, g^{(d)}$; we shall construct a new circuit $C'$ such that $g^{(i)}$ computes the degree $i$ homogeneous component of the polynomial computed at $g$. 
If $g$ has children $h_1$ and $h_2$, then $C'$ would have the following connections depending on the type of the gate $g$:
\begin{eqnarray*}
\text{$g = h_1 + h_2$}\quad\implies\quad g^{(i)} &=& h_1^{(i)} \spaced{+} h_2^{(i)}\quad\text{for all $i$}\\
\text{$g = h_1 \times h_2$}\quad\implies\quad g^{(i)} &=& \sum_{j=0}^i h_1^{(j)} h_2^{(i-j)}\quad\text{for all $i$}
\end{eqnarray*}
It is easy to check that the size of the circuit $C'$ is at most $O(sd^2)$, and computes all the homogeneous components of $f$. 
\end{proof}

Thus, for arithmetic circuits, we can assume without much loss of generality that we are working with a homogeneous circuit. 

\begin{remark*}
For the class of arithmetic formulas, it is not clear if we can homogeneous without loss of generality. 
If we were to apply the above lemma to an arbitrary arithmetic formula, the resulting object is a homogeneous circuit and not a formula. 
It is unclear if any formula can be homogenized without loss of generality. 
The same is the case even for circuits of a fixed depth, as the above construction doubles the depth of the circuit. 

However, the class of ABPs can also be assumed to be homogeneous without loss of generality. 
We leave this as an exercise. 
\end{remark*}

\begin{exercise}
Prove a similar homogenization lemma for algebraic branching programs. 
\end{exercise}

\section{Interpolation}

A very useful technique in arithmetic complexity is interpolation. Suppose we have a circuit $C$ that computes a polynomial $P(x_1,\ldots, x_n, y)$ and say $\deg_y P = d$. We can interpret the polynomial $P$ as an element of $(\F[\vecx])[y]$, that is we can think of $P$ as a univariate polynomial in $y$ with coefficients being elements of $\F[\vecx]$. 
\[
P(x_1,\ldots, x_n, y) \spaced{=} P_0(x_1,\ldots, x_n) + P_1(x_1,\ldots, x_n) y + \cdots + P_d(x_1,\ldots, x_n) y^d
\]
Given that we can compute $P$ by a circuit $C$, can we also compute one of the coefficients $P_i(x_1,\ldots, x_n)$ by a small circuit? The answer is `Yes' as long as we are working over a large enough field. 

\begin{lemma}[Interpolation]\label{lem:interpolation}
Let $P(x_1,\ldots, x_n, y)$ be a polynomial with $\deg_y P \leq d$. For any set of distinct scalars $\alpha_0, \ldots, \alpha_{d} \in \F$ and for every $i \in \set{0,\ldots, d}$,  each of the coefficient polynomials $P_i(x_1,\ldots, x_n)$ defined above can be expressed as a linear combination of 
\[
\set{P(x_1,\ldots, x_n, \alpha_0), \ldots, P(x_1,\ldots, x_n,\alpha_d)}.
\]
Thus in particular, if $P$ is computed by a size $s$ circuit, then each $P_i$ is computable  by a size $s(d+1)$ circuit. Furthermore, if $P$ is computable by a size $s$ circuit from some class $\mathcal{C}$, then each $P_i$ is computable by a size $s(d+1)$ circuit from the class $\Sigma \mathcal{C}$ (which are circuits with a $+$ gate as a root with $\mathcal{C}$-circuits as its children).
\end{lemma}
\begin{proof}
For this proof we shall use $P(\alpha_i)$ as a short-hand for $P(x_1,\ldots, x_n, \alpha_i)$. Then we have the following matrix identity
\[
\insquare{\begin{array}{cccc}
1 & \alpha_0 & \cdots & \alpha_0^d\\
1 & \alpha_1 & \cdots & \alpha_1^d\\
\vdots & \vdots & \ddots & \vdots\\
1 & \alpha_d & \cdots & \alpha_d^d
\end{array}} \insquare{\begin{array}{c}
P_0\\
P_1 \\
\vdots\\
P_d
\end{array}} \spaced{=}
\insquare{\begin{array}{c}
P(\alpha_0)\\
P(\alpha_1) \\
\vdots\\
P(\alpha_d)
\end{array}}
\]
Observe that the matrix on the left is a Vandermonde matrix, and hence is invertible. This implies that there is some matrix $\inparen{\!\inparen{\beta_{ij}}\!}$, its inverse, such that
\[
\insquare{\begin{array}{c}
P_0\\
P_1 \\
\vdots\\
P_d
\end{array}} \spaced{=}
\insquare{\begin{array}{cccc}
\beta_{00} & \beta_{01} & \cdots & \beta_{0d}\\
\beta_{10} & \beta_{11} & \cdots & \beta_{1d}\\
\vdots & \vdots & \ddots & \vdots\\
\beta_{d0} & \beta_{d1} & \cdots & \beta_{dd}
\end{array}} 
\insquare{\begin{array}{c}
P(\alpha_0)\\
P(\alpha_1) \\
\vdots\\
P(\alpha_d)
\end{array}}
\]
In particular, for each $i \in \set{0,\ldots, d}$,
\[
P_i(x_1,\ldots, x_n) \spaced{=} \beta_0 P(x_1,\ldots, x_n, \alpha_0) + \cdots + \beta_d P(x_1,\ldots, x_n, \alpha_d). 
\]
The size bounds and structure claimed is clear from the above expression.
\end{proof}

The use of interpolation is one of the most important methods when dealing with arithmetic circuits. We would many applications of it later in this article. We just give a couple of concrete applications for now. The first is the computation of partial derivatives in a single variable. We leave this as an easy exercise. 

\begin{exercise}[Computing partial derivatives]
Suppose $P(x_1,\ldots, x_n,y)$ be a polynomial computed by a size $s$ circuit with $\deg_y P \leq d$. Show that for any $i$, the partial derivative
$\frac{\partial^i P}{\partial y^i}$ can be computed by a circuit of size at most $s(d+1)$. 

What about computing mixed partials such as $\frac{\partial^n P}{\partial x_1 \cdots \partial x_n}$?
\end{exercise}

\noindent 
Another important application is the task of computing homogeneous components. 

\subsection{Computing homogeneous components}

Suppose we have a size $s$ circuit $C$ that is computing a polynomial $P$ of degree $d$.
We have already seen in \autoref{lem:homogenization} that we can construct each of the homogeneous components of $P$ by a circuit $C'$ of size at most $sd^2$.
Unfortunately, the construction described in \autoref{sec:homogenization} results in a circuit $C'$ even if we started off with a formula computing $P$.
Furthermore, if $P$ was computable by a constant depth circuit, the construction results in an non-constant depth circuit.
However, the construction in \autoref{lem:homogenization} not just computes the homogeneous components of $P$ but it computes them via a \emph{homogeneous} circuit. If the goal is to just compute the homogeneous components (by a possibly non-homogeneous circuit), one can use interpolation. Furthermore, this method would preserve constant depth as well.

Let us use $\Hom_i(P)$ to denote the homogeneous part of degree $i$, and let $d = \deg(P)$.
\[
P(x_1,\ldots, x_n) \spaced{=} \Hom_0(P) + \cdots + \Hom_d(P)
\]
Let $y$ be a new variable and consider the polynomial $P'(x_1,\ldots, x_n, y) := P(yx_1, \ldots, y x_n)$. Then observe that
\[
P'(x_1,\ldots, x_n) \spaced{=} \Hom_0(P) + y\Hom_1(P) + \cdots + y^d \Hom_d(P).
\]
In other words, each homogeneous component of $P$ is a coefficient of some $y^i$ of $P'(x_1,\ldots, x_n,y)$. Thus, \autoref{lem:interpolation} tells us that we can use a linear combination of evaluations of $P'$ to compute these coefficients. 

We summarize this as a lemma. 

\begin{lemma}[Homogeneous component computations]\label{lem:computing-hom-components}
Let $P(x_1,\ldots, x_n)$ be a polynomial with $\deg P = d$. For any set of distinct scalars $\alpha_0, \ldots, \alpha_{d} \in \F$ and for every $i \in \set{0,\ldots, d}$,  each of the homogeneous components $\Hom_i(P)$ defined above can be expressed as a linear combination of 
\[
\set{P(\alpha_0x_1,\ldots, \alpha_0x_n), \ldots, P(\alpha_dx_1,\ldots, \alpha_dx_n)}.
\]
Thus in particular, if $P$ is computed by a size $s$ circuit, then each $\Hom_i(P)$ is computable  by a size $s(d+1)$ circuit. Furthermore, if $P$ is computable by a size $s$ circuit from some class $\mathcal{C}$, then each $\Hom_i(P)$ is computable by a size $s(d+1)$ circuit from the class $\Sigma \mathcal{C}$ (which are circuits with a $+$ gate as a root with $\mathcal{C}$-circuits as its children).
\end{lemma}

It is important to note that for interpolation, we need to be able to find sufficiently many distinct elements in the base field, and this is possible only if the base field is large enough. 

\begin{exercise}\label{ex:ben-or-trick}
Assume the underlying field is large enough. Find a polynomial sized constant depth circuit that computes $\ESym_d$, the elementary symmetric polynomial of degree $d$, defined as
\[
\ESym_{d}(x_1,\ldots, x_n) \spaced{=} \sum_{1\leq i_1 < i_2 < \cdots < i_d \leq n} x_{i_1} \cdots x_{i_d}.
\]
\end{exercise}

\begin{exercise}
Assume the underlying field is large enough. Find a polynomial sized constant depth circuit that computes \emph{complete homogeneous symmetric polynomial} defined as
\[
\sigma_d(x_1,\ldots, x_n) \spaced{=} \sum_{m \in \set{\text{deg. $d$ monomials}}} m.
\]
\textcolor{gray}{For example, $\sigma_2(x_1,x_2,x_3) = x_1^2 + x_2^2 + x_3^2 + x_1x_2 + x_2x_3 + x_1x_3$.}
\end{exercise}

%%% Local Variables: 
%%% mode: latex
%%% TeX-master: "fancymain"
%%% End: 
